%% 51장 모멘트, 생성함수, 극한 정리
\chapter{Moments, Generating Functions, and Limit Theorems}
\label{ch:limit-theorems}

\abstract{분포를 몇 개의 수로 요약하는 모멘트와, 분포 전체를 하나의 함수로 담는 생성함수를 다룬다. 적률생성함수와 누적량은 독립 확률변수의 합을 곱과 덧셈으로 바꾸고 특성함수는 모든 분포에 대해 존재하며 분포를 결정한다. 이 도구로 통계학이 기대는 두 극한 정리인 큰 수의 법칙과 중심극한정리를 증명한다. 약한 법칙은 Chebyshev 부등식으로, 강한 법칙은 4차 모멘트를 가정하고 Borel--Cantelli 보조정리로, 중심극한정리는 특성함수의 Taylor 전개와 Lévy의 연속성 정리로 보인다. 지수 분포의 합에서는 정규 근사의 오차가 점근적으로 $0.1330/\sqrt n$으로 줄어드는데 그 상수가 Berry--Esseen 경계가 아니라 Edgeworth 전개의 첫 항에서 나옴을 확인한다. 끝으로 표본 모멘트를 binary64로 계산하는 문제를 다룬다. 평균이 큰 자료에서 교과서의 한 번 훑기 공식은 분산을 완전히 잃고 두 번 훑기도 3차 모멘트에서 열 자리 가까이를 잃어 여섯 자리 남짓만 남는데 stardust는 평균을 두 수의 합으로 들고 다녀 이것을 막는다.}

\section{Moments and Cumulants}
\label{sec:moments}

확률변수 $X$의 $k$차 \term{모멘트}{moment} $\mu_k'=E[X^k]$와 평균 $\mu=E[X]$ 둘레의 \term{중심 모멘트}{central moment} $\mu_k=E[(X-\mu)^k]$는 분포의 모양을 요약한다. $\mu_2=\sigma^2$이 분산이고
\[
\gamma_1=\frac{\mu_3}{\sigma^3},\qquad\gamma_2=\frac{\mu_4}{\sigma^4}-3
\]
을 각각 \term{왜도}{skewness}와 \term{초과 첨도}{excess kurtosis}라 하며 정규분포에서는 둘 다 $0$이다. 왜도는 꼬리의 비대칭을, 초과 첨도는 정규분포에 견준 꼬리의 무게를 잰다. 지수 분포는 $\gamma_1=2$, $\gamma_2=6$이다.

모멘트보다 다루기 쉬운 양으로 누적량이 있다. $X$의 \term{적률생성함수}{moment generating function}를 $M(t)=E[e^{tX}]$로 정의하자. $0$의 근방에서 $M$이 유한하면 거듭제곱 급수 $M(t)=\sum_k\mu_k't^k/k!$로 모멘트가 모두 나오고, $K(t)=\log M(t)$의 Taylor 계수 $K(t)=\sum_k\kappa_kt^k/k!$를 \term{누적량}{cumulant}이라 한다. $\kappa_1=\mu$, $\kappa_2=\sigma^2$, $\kappa_3=\mu_3$, $\kappa_4=\mu_4-3\sigma^4$이므로 $\gamma_1=\kappa_3/\kappa_2^{3/2}$, $\gamma_2=\kappa_4/\kappa_2^2$이다. 누적량은 독립 확률변수의 합에서 더해지므로 편리하다. $X$와 $Y$가 독립이면 $E[e^{t(X+Y)}]=E[e^{tX}]E[e^{tY}]$이므로 $K_{X+Y}=K_X+K_Y$이고 모든 차수의 누적량이 더해진다. 정규분포 $N(\mu,\sigma^2)$은 $K(t)=\mu t+\sigma^2t^2/2$여서 3차 이상의 누적량이 모두 $0$이다.

이 덧셈 성질에서 중심극한정리의 모습이 이미 보인다. 같은 분포를 따르는 독립 확률변수 $n$개의 합을 표준화한 $S_n^\ast=(S_n-n\mu)/(\sigma\sqrt n)$의 $k$차 누적량은 $k\ge2$에서 $n\kappa_k/(\sigma\sqrt n)^k=n^{1-k/2}\kappa_k/\sigma^k$이고 1차 누적량은 $0$이다. 따라서 $k\ge3$의 누적량이 모두 $0$으로 가고, 남는 것은 평균 $0$, 분산 $1$, 곧 표준 정규분포의 누적량이다. 왜도는 $\gamma_1/\sqrt n$, 초과 첨도는 $\gamma_2/n$의 비율로 줄어든다.

\section{Characteristic Functions}
\label{sec:characteristic-function}

적률생성함수는 꼬리가 무거운 분포에서 $0$ 밖의 모든 $t$에 대해 무한대일 수 있다. 로그정규분포는 모든 차수의 모멘트가 있지만 적률생성함수가 $t>0$에서 무한대이고, 같은 모멘트를 가진 다른 분포가 있다. 그래서 일반적인 도구로는
\[
\varphi(t)=E[e^{itX}]=E[\cos tX]+iE[\sin tX],\qquad t\in\mathbb R
\]
인 \term{특성함수}{characteristic function}를 쓴다. $|e^{itX}|=1$이므로 특성함수는 모든 분포에 대해 존재하고 $|\varphi|\le1$이며 $\varphi(0)=1$이고 지배 수렴 정리로 고르게 연속이다. 특성함수는 37장의 Fourier 변환에서 지수의 부호만 바꾼 것, 곧 $\varphi(t)=\widehat{P_X}(-t)$이다.

특성함수의 중요한 성질은 셋이다.
\begin{enumerate}
\item 독립 확률변수의 합의 특성함수는 특성함수의 곱이다.
\item 특성함수는 분포를 결정한다. 반전 공식
\[
P(a<X<b)+\frac12P(X\in\{a,b\})=\lim_{T\to\infty}\frac1{2\pi}\int_{-T}^T\frac{e^{-ita}-e^{-itb}}{it}\varphi(t)\,dt
\]
이 특성함수에서 분포를 되찾는다\cite{billingsley1995}.
\item (Lévy의 연속성 정리) 분포들의 열 $P_n$이 $P$로 약하게 수렴하는 것은 특성함수가 모든 $t$에서 $\varphi_n(t)\to\varphi(t)$인 것과 동치이다\cite{billingsley1995}.
\end{enumerate}
여기서 \term{약한 수렴}{weak convergence}, 곧 \term{분포 수렴}{convergence in distribution}은 극한의 누적분포함수 $F$가 연속인 모든 점에서 $F_n(x)\to F(x)$인 것이다.

$E|X|^k<\infty$이면 특성함수는 $k$번 미분 가능하고 $\varphi^{(j)}(0)=i^jE[X^j]$($j\le k$)이다. 특히 $E[X]=0$, $E[X^2]=\sigma^2$이면
\begin{equation}
\varphi(t)=1-\frac{\sigma^2t^2}2+o(t^2)\qquad(t\to0)
\label{eq:cf-taylor}
\end{equation}
이다. 이 전개는 2차 모멘트만 있으면 성립한다. $e^{ix}$의 Taylor 나머지에 관한 부등식 $|e^{ix}-1-ix+x^2/2|\le\min(|x|^3/6,\,x^2)$을 $x=tX$에 쓰고 지배 수렴 정리를 쓰면 된다.

표준 정규분포의 특성함수는 $e^{-t^2/2}$이고, 표준 코시 분포의 특성함수는 $e^{-|t|}$이다. 그러면 독립인 코시 확률변수 $n$개의 평균의 특성함수는 $(e^{-|t|/n})^n=e^{-|t|}$로 하나의 코시와 같아서 평균을 아무리 많이 내도 흩어짐이 줄지 않는다. 평균이 없으면 강한 법칙은 반드시 깨지고(약한 법칙은 꼬리에 따라 성립하기도 한다), 코시 분포는 표본평균이 전혀 모이지 않는 극단적인 예이다.

\section{Laws of Large Numbers}
\label{sec:lln}

같은 분포를 따르는 독립 확률변수 $X_1,X_2,\dots$의 표본평균을 $\bar X_n=S_n/n$이라 하자.

\begin{theorem}[약한 법칙]
$E[X_1^2]<\infty$이면 모든 $\epsilon>0$에 대해 $P(|\bar X_n-\mu|>\epsilon)\to0$이다.
\end{theorem}

\begin{proof}
독립이므로 $\operatorname{Var}\bar X_n=\sigma^2/n$이고, Chebyshev 부등식(49장)으로 $P(|\bar X_n-\mu|>\epsilon)\le\sigma^2/(n\epsilon^2)\to0$이다.
\end{proof}

\begin{theorem}[강한 법칙, 4차 모멘트]
\label{thm:slln4}
$E[X_1^4]<\infty$이면 $\bar X_n\to\mu$가 거의 확실하다.
\end{theorem}

\begin{proof}
$\mu=0$이라 해도 된다. $E[S_n^4]=\sum_{i,j,k,l}E[X_iX_jX_kX_l]$에서 첨자 하나가 홀로 나타나는 항은 독립성과 $E[X_i]=0$으로 $0$이다. 남는 것은 네 첨자가 같은 $n$개 항과 두 쌍으로 나뉘는 $3n(n-1)$개 항이므로 $E[S_n^4]=nE[X^4]+3n(n-1)\sigma^4\le Cn^2$이다. Markov 부등식으로
\[
P(|\bar X_n|>\epsilon)=P(S_n^4>\epsilon^4n^4)\le\frac{Cn^2}{\epsilon^4n^4}=\frac C{\epsilon^4n^2}
\]
이고 이 상계를 $n$에 대해 더한 합은 유한하다. 49장의 Borel--Cantelli 보조정리로 $|\bar X_n|>\epsilon$인 일은 거의 확실히 유한 번만 일어난다. $\epsilon=1/m$, $m=1,2,\dots$의 가산 개 사건을 모으면 $\bar X_n\to0$이 거의 확실하다.
\end{proof}

강한 법칙은 사실 $E|X_1|<\infty$만으로 성립하고 Etemadi는 이를 쌍마다 독립인 경우까지 넓혀 짧게 증명했다\cite{etemadi1981}. 몬테카를로 방법(60장)은 이 법칙에 기댄다. 표본평균이 참평균으로 거의 확실히 다가가므로 충분히 많이 뽑으면 기댓값을 원하는 만큼 정확히 어림할 수 있다. 다가가는 속도는 다음 정리에서 드러난다.

\section{The Central Limit Theorem}
\label{sec:clt}

\begin{theorem}[Lindeberg--Lévy]
\label{thm:clt}
$X_i$들이 독립이고 같은 분포를 따르며 $E[X_1]=\mu$, $0<\operatorname{Var}X_1=\sigma^2<\infty$이면
\[
\frac{S_n-n\mu}{\sigma\sqrt n}\ \xrightarrow{\ d\ }\ N(0,1)
\]
이다.
\end{theorem}

\begin{proof}
$\mu=0$, $\sigma=1$이라 해도 된다. 표준화한 합의 특성함수는 $\varphi(t/\sqrt n)^n$이다. \eqref{eq:cf-taylor}로 $\varphi(t/\sqrt n)=1-t^2/(2n)+o(1/n)$이다. 절댓값이 $1$ 이하인 복소수 $a_j$, $b_j$에 대해 $|\prod a_j-\prod b_j|\le\sum|a_j-b_j|$이므로(항을 하나씩 바꾸어 망원 합으로 쓴다)
\[
\Bigl|\varphi(t/\sqrt n)^n-\Bigl(1-\frac{t^2}{2n}\Bigr)^n\Bigr|\le n\cdot o(1/n)\to0
\]
이다($n$이 크면 $|1-t^2/(2n)|\le1$이다). $(1-t^2/(2n))^n\to e^{-t^2/2}$이므로 특성함수가 표준 정규분포의 특성함수로 수렴하고 Lévy의 연속성 정리로 분포가 수렴한다.
\end{proof}

Lindeberg는 1922년에 분포가 같지 않더라도 각 항이 합의 분산에 견주어 무시할 만하다는 조건 아래 같은 결론이 성립함을 증명했다\cite{lindeberg1922}. 분포마다 평균 $0$, 분산 $\sigma_i^2$이고 $s_n^2=\sum_{i\le n}\sigma_i^2$이라 할 때 그 조건은 모든 $\epsilon>0$에 대해
\[
\frac1{s_n^2}\sum_{i=1}^nE\bigl[X_i^2\mathbf 1\{|X_i|>\epsilon s_n\}\bigr]\to0
\]
이다. 회귀 계수의 추정량처럼 서로 다른 무게의 합에 정규 근사를 쓸 때 이 형태가 필요하다.

\subsection{Rate of Convergence}

중심극한정리 자체는 수렴이 얼마나 빠른지 알려 주지 않는다. Berry\cite{berry1941}와 Esseen은 3차 절대 모멘트 $\rho=E|X_1-\mu|^3$이 유한하면
\[
\sup_x\bigl|F_n(x)-\Phi(x)\bigr|\le\frac{C\rho}{\sigma^3\sqrt n}
\]
임을 보였다. $F_n$은 표준화한 합의 누적분포함수이다. 상수 $C$의 가장 좋은 값은 지금도 연구되고 있으며, 같은 분포의 경우 $C\le0.4748$이 알려져 있다\cite{shevtsova2011}.

이 경계가 얼마나 날카로운지 지수 분포로 확인해 보자. 평균 $1$인 지수 확률변수 $n$개의 합은 모양 $n$인 감마 분포를 따르므로 $F_n$을 정확히 계산할 수 있다. $\rho=E|X-1|^3=12/e-2=2.4146$이고 경계는 $1.146/\sqrt n$이다. 실제 최대 오차는 $n=1,10,100,1000,10^4$에서 $0.159$, $0.0421$, $0.0133$, $0.00421$, $0.00133$이다. $\sqrt n$을 곱하면 $0.159$, $0.133$, $0.1330$, $0.1330$, $0.1330$으로, $n$이 커지면 $1/\sqrt n$의 비율로 줄되 경계보다 약 $8.6$배 작다. 상수 $0.1330$은 \term{Edgeworth 전개}{Edgeworth expansion}의 첫 보정항에서 나온다. 분포가 격자 분포가 아니고(Cramér 조건) 4차 모멘트가 있으면 누적량을 써서 $F_n$을 전개할 수 있고
\[
F_n(x)=\Phi(x)-\frac{\gamma_1}{6\sqrt n}(x^2-1)\phi(x)+O(1/n)
\]
이고, 보정항의 크기는 $x=0$에서 가장 커서 $\gamma_1\phi(0)/6=2\times0.39894/6=0.13298$이다. 이 식을 보면 오차는 주로 왜도에서 오고 대칭이고 격자가 아닌 분포에서는 정규 근사가 $1/n$의 비율로 더 빨리 좋아진다. 동전 던지기 같은 격자 분포는 대칭이어도 $F_n$의 계단 때문에 오차가 $1/\sqrt n$의 비율에 머문다.

\begin{aside}[정규 분포와 오차의 법칙]
정규 분포는 de Moivre가 1733년에 이항 분포의 근사로 처음 얻었고 Laplace가 일반적인 합으로 넓혔으며 Gauss가 측정 오차의 분포로 썼다. ``중심극한정리''라는 이름은 확률론에서 중심이 되는 극한 정리라는 뜻으로 Pólya가 1920년에 붙였다\cite{polya1920}. Lyapunov가 1901년에 특성함수로 처음 엄밀한 일반 증명을 냈고, Lindeberg의 1922년 조건과 Lévy의 연속성 정리로 오늘날의 형태가 갖추어졌다.
\end{aside}

\section{Computing Sample Moments}
\label{sec:sample-moments}

표본 $x_1,\dots,x_n$의 분산을 $n^{-1}\sum x_i^2-\bar x^2$로 계산하는 공식은 수학적으로 옳지만 두 항이 거의 같은 큰 수이면 9장의 소거가 결과를 삼키므로 binary64에서는 쓸 수 없다. 평균 $10^8$, 표준편차 $1$인 정규 표본 $10^4$개를 만들어 정확한 유리수 산술로 분산 $0.996994\ldots$을 구하고 여러 방법과 견주면 다음과 같다.
\begin{itemize}
\item 한 번 훑기 공식은 numpy의 \code{mean}으로 계산하면 $0$, 단순한 고리로 더하면 $90$이 나온다. $x_i^2\approx10^{16}$의 반올림 간격이 $2$이므로 차이 $1$은 남지 않는다.
\item 평균을 먼저 구하고 편차의 제곱을 평균하는 두 번 훑기는 상대오차 $1.2\times10^{-16}$이다.
\item 평균을 한 원소씩 갱신하는 Welford의 방법\cite{welford1962}은 상대오차 $2.9\times10^{-9}$이다. 평균 $10^8$ 근처에서 갱신할 때마다 반올림이 쌓인다.
\end{itemize}
두 번 훑기도 완전하지는 않다. 반올림한 평균 $\hat m$은 정확한 표본평균 $\bar x$와 $\delta=\hat m-\bar x$만큼 다르고, 그 크기는 평균을 정확히 반올림했다면 $10^8$의 반올림 간격의 절반인 $7.5\times10^{-9}$까지일 수 있다. 이 예에서는 $\delta=5.9\times10^{-9}$이다. 편차를 $\hat m$에서 재면 2차 모멘트는 $\delta^2$만큼만 틀리지만 3차 모멘트는
\[
\frac1n\sum(x_i-\hat m)^3=m_3-3\delta m_2-\delta^3
\]
이어서 $3\delta m_2\approx1.8\times10^{-8}$만큼 틀린다. 정규 표본의 $m_3$은 작으므로($-0.00566$) 상대오차가 $3.1\times10^{-6}$에 이른다. 평균의 오차까지 계산해 보정하면 이 오차를 없앨 수 있다. 편차 $d_i=x_i-\hat m$의 평균 $S_1=\bar d$가 바로 $-\delta$이므로, $S_k=n^{-1}\sum d_i^k$에서
\[
m_2=S_2-S_1^2,\qquad m_3=S_3-3S_1S_2+2S_1^3,\qquad m_4=S_4-4S_1S_3+6S_1^2S_2-3S_1^4
\]
로 고치면 같은 두 번 훑기로 $m_3$의 상대오차가 $6.4\times10^{-16}$이 된다. Chan, Golub, LeVeque는 이런 알고리즘들의 오차를 자료의 조건수 $\kappa=\sqrt{1+n\bar x^2/S}$($S=\sum(x_i-\bar x)^2$, 9장)로 분석했다\cite{chan1983variance}. 이 예의 조건수는 $10^8$이다.

\begin{instardust}[\code{docs/design/stats.md} §5, T69]
\code{sd.stats}의 기술통계(\code{describe}, \code{skew}, \code{kurtosis}, \code{moment}, \code{zscore} 등)는 이 절의 방법으로 모멘트를 계산한다. 평균을 $\text{hi}+\text{lo}$의 두 수로 들어 첫 훑기의 $\text{hi}=\text{합}/n$은 어림이고 둘째 훑기에서 $d=x-\text{hi}$의 $S_k$를 구한 뒤 $\text{lo}=S_1$로 위의 식처럼 $m_2,m_3,m_4$를 고치며 세 번째로 훑지는 않는다. 합은 128원소 블록마다 여덟 개의 사슬로 나누어 더하는데, 결과가 오차를 키우는 합(3차 모멘트의 $d$와 $d^3$처럼)은 오차를 지니는 TwoSum 사슬로 더하고, 블록의 결과는 8192원소 블록 안에서 Neumaier의 보상 합으로, 큰 블록들은 그 합과 보정을 차례로 더한다. 자르는 위치가 원소의 번호로만 정해지므로 스레드 수와 상관없이 결과가 비트까지 같다. 첨도는 scipy와 같이 초과 첨도가 기본이고, \code{bias=False}이면 표본 크기에 따른 수정 계수를 곱한다. 상수 표본의 왜도와 첨도는 NaN이다.
\end{instardust}

\section{Concluding Remarks}

모멘트와 누적량은 분포를 수로 요약하고, 생성함수는 독립 확률변수의 합을 곱으로 바꾼다. 특성함수는 언제나 존재하고 분포를 결정하며 그 수렴이 분포의 수렴과 같으므로 극한 정리의 자연스러운 도구가 된다. 큰 수의 법칙에 따르면 표본평균은 참평균으로 가고 중심극한정리에 따르면 그 오차는 근사적으로 표준편차 $\sigma/\sqrt n$인 정규 분포를 따른다. Berry--Esseen 경계와 Edgeworth 전개로는 그 근사가 얼마나 좋은지를 가늠할 수 있는데 실제 오차는 왜도가 정하는 Edgeworth의 첫 항으로 잘 설명된다. 모멘트를 binary64로 계산할 때에도 조심할 점이 있어서 평균이 큰 자료에서는 평균 자체의 반올림까지 보정해야 3차 이상의 모멘트가 제 자릿수를 지킨다. 다음 장에서는 여러 확률변수를 함께 다루는 다변량 분포로 넘어간다.

극한 정리는 Billingsley\cite{billingsley1995}와 Durrett\cite{durrett2019}의 책에, 표본 분산의 알고리즘은 Chan, Golub, LeVeque의 논문\cite{chan1983variance}과 Higham의 책\cite{higham2002accuracy}에 자세하다.

\chapterreferences
